\section{A constrained reference is not a Gaussian probability law}
There are two different uses of the same exponential expression. One may
define a geometric score $V(x)$ and note that $\exp[-V(x)]$ has a
Gaussian shape. Or one may define a probability measure proportional to
that expression on a specified domain, with a specified base measure and
normalizing constant. The second is an additional model. It does not follow
from the first and does not describe an actor process until a connection
to that process is established.

Conservation makes the distinction concrete. In two cells with total mass
twenty, choosing $x_1$ determines $x_2=20-x_1$. They cannot be independent
random coordinates on that mass line. A diagonal Hessian still gives a
separable formula on the ambient space, but the feasible states are
coupled by the constraint. A simulation histogram resembling a bell would
also not prove independent Gaussian coordinates.

The same care applies to the scale $\sigma_i$. In this book it is a
declared width in physical stock units. Calling it a standard deviation
would require a specified probability distribution and calibration. Calling
it a measurement error would require a sensor-error model. Neither role
is obtained simply because statistical notation often uses the same letter.

\section{The reference and the directions the world can use}
In an unconstrained ambient space, a vanishing gradient characterizes the
minimum of this strictly convex quadratic. A conserved world moves only
in directions whose entries sum to zero. If every allowed directional
derivative is zero, the state need not have zero ambient gradient in a
more general incompatible-reference problem. Equal marginals can make
all conservative first-order transfers neutral while the individual
marginals are nonzero.

The active compatible-reference model avoids that particular ambiguity at
its global minimum, but graph restrictions can create another. If the
graph has disconnected components, no internal action can change a
component's total mass. A reference compatible with global mass may still
be incompatible with a component's available mass. Thus global compatibility
is not a reachability theorem even for an otherwise continuous action set.

One can see this without running a process. Divide four cells into two
disconnected pairs. If the first pair holds twelve units but its two
references sum to twenty, no transfer confined to that pair can meet both
references. Movement elsewhere does not repair the missing route. A global
potential may show a persistent deviation that is structurally impossible
to remove under the declared topology.

\section{Why a small potential is an incomplete functional statement}
Because the potential sums standardized squared deviations, a small value
does bound each represented standardized deviation. In fact,
\begin{equation}
 \frac{|x_i-x_i^*|}{\sigma_i}\leq\sqrt{2V(x)}
 \quad\hbox{for every represented cell }i.
\end{equation}
The proof uses nonnegativity: an individual summand cannot exceed the
whole sum. This is a useful geometric guarantee. It concerns represented
coordinates and declared scales only.

It does not establish the adequacy of the reference for health, the
quality of the resource, access for an unrepresented household or the
absence of damage outside the system boundary. Those concerns are not
counterexamples to the inequality. They are reasons to avoid assigning
the inequality a meaning it never had. Functional calibration and
representational completeness require additional evidence and choices.

Conversely, a large potential does not prove that one actor caused the
departure or owes compensation for it. Initial conditions and external
forcing can create deviation before the actor chooses anything. The
following chapters therefore bind value to a specific transition and
baseline, rather than treating the state score as a debt automatically
assigned to whoever happens to be present.
